\section{Conclusion}
\label{sec:conclusion}

\subsection{The Binding Constraint Diagnosis}

This paper has argued that port connector segments are the binding capacity constraint on US import and export freight throughput for at least five major ports---and that this constraint is systematically invisible in the standard T1 corridor analysis and freight bottleneck ranking methodologies. Seven of ten port connectors operate at V/C $> 1.2$ at peak hour. The worst---I-710 at Long Beach---operates at V/C 1.52 and accumulates an estimated 8.3 million truck-delay-hours per year, four times the top-ranked ATRI bottleneck, while not appearing on the ATRI bottleneck list at all.

The invisibility is not an oversight. It is structural: the HPMS and ATRI methodologies were designed to analyze corridor-scale infrastructure, and port connectors are too short to register as corridors. The measurement framework produces a systematic false negative for the constraint class that matters most to international trade competitiveness. The paper's primary methodological contribution is the adjusted peak-hour V/C methodology (Equation~\ref{eq:vc}) that corrects for the drayage demand concentration that makes port connectors uniquely severe in their peak-hour performance relative to their 24-hour average.

\subsection{Summary of Findings}

The four research questions are answered as follows:

\begin{enumerate}
  \item \textbf{Congestion baseline}: Seven of ten port connectors exceed V/C 1.2 at peak hour under the adjusted methodology. The worst connector (I-710) exceeds V/C 1.5 and generates approximately \$1.86 billion per year in direct truck delay cost.

  \item \textbf{Failure mode taxonomy}: Four failure modes are identified. Capacity failure (Mode A) is the most common, affecting all ten connectors to varying degrees. Classification failure (Mode B) affects Seattle/Tacoma, where SR-509 carries near-interstate volumes on a sub-interstate geometric standard. Resilience failure (Mode C) affects LA/Long Beach and Norfolk, which have no viable alternate connector route. Jurisdictional fragmentation (Mode D) affects six of ten connectors and is the primary cause of investment delay at Oakland, NY/NJ, Charleston, Miami, and Baltimore.

  \item \textbf{USMCA border case}: The I-35 connector at Laredo generates an estimated \$10.8 billion per year in total delay cost from 16,000 trucks per day with 2--4 hour wait times. A dedicated truck lane on the 8-mile connector costs \$280--360 million and produces a benefit-cost ratio of 14:1 to 108:1 depending on the benefit capture assumption---the highest-BCR freight investment identified in the ROUTE corpus.

  \item \textbf{Investment case}: A portfolio of port connector upgrades totaling \$20--50 billion increases national import/export freight throughput by approximately 15\% at a cost per throughput-unit 40--60\% lower than equivalent T1 managed-lane expansion. The LA/Long Beach I-710 improvement dominates the portfolio at \$3.2--4.8 billion and 9 of the 15 percentage points of throughput improvement.
\end{enumerate}

\subsection{The Port Connector Standard}

The central policy recommendation is the establishment of a ``port connector'' classification in the National Highway System under 23 USC 103(c). The classification would:

\begin{itemize}
  \item Apply to the primary approach road from any top-25 US port gate to the first T1 or T2 interstate interchange.
  \item Mandate automatic NHS designation for classified segments, enabling NHPP and INFRA funding eligibility.
  \item Establish a design standard based on drayage peak demand (95th-percentile peak hour at V/C $\leq 0.85$) rather than 24-hour AADT.
  \item Require real-time V/C monitoring integrated with port gate appointment systems.
  \item Create a coordinating authority with cross-jurisdictional investment authority, resolving the fragmentation that delays investment at six of ten analyzed connectors.
\end{itemize}

This recommendation does not require new authorization beyond what is available under IIJA \citep{IIJA2021}: the NHPP, INFRA, and Port Infrastructure Development Program (PIDP) collectively authorized \$17.5 billion for port-adjacent freight investment. What is missing is the classification mechanism that directs that funding to the right segments.

\subsection{Savannah: Act Now}

The Savannah finding merits special emphasis because it has a decision deadline. At a 12\% CAGR, Savannah's port connector will constrain throughput below the expanded terminal capacity by approximately 2028. Investment lead time is 7--9 years. A decision to invest made in 2026 produces infrastructure relief in 2033--2035; a decision delayed to 2028 produces relief in 2035--2037. The GPA's \$2.5 billion terminal expansion investment, completed in 2023, will be throttled by a \$890 million--\$1.23 billion connector problem if the decision is deferred.

The Savannah case is generalizable: port connector investment must be initiated before the constraint binds, not after, because the lead time exceeds the growth window. A ``port connector early warning'' program---monitoring V/C trajectory at all classified connectors and triggering project initiation when V/C exceeds 1.0 at peak---would prevent the Savannah scenario from being repeated at Charleston (V/C 1.09 and growing) or Baltimore (V/C 1.12 and growing) in the decade that follows.

\subsection{Limitations and Next Steps}

The V/C estimates in Table~\ref{tab:vc} are computed from HPMS 2021 data \citep{HPMS2018} adjusted for port drayage peak demand profiles. Three connectors required AADT estimation from BTS port gate data rather than direct HPMS measurement, introducing uncertainty of approximately $\pm$15\% in those cases. The throughput improvement estimates (15\% national aggregate) are sensitive to the assumed port growth rates: at 6\% CAGR (half the recent Savannah rate), the binding date at Savannah is 2035 rather than 2028, reducing the urgency premium but not the BCR of the investment.

Three extensions would strengthen the analysis: (1) direct loop-detector V/C measurement on all ten connectors at 15-minute resolution (available from Caltrans PeMS for I-710 but not standardized nationally); (2) vessel demurrage cost quantification (the cost of ships waiting at anchor because truck congestion slows terminal throughput) which the current analysis excludes; and (3) a network-level analysis of connector investment sequencing that captures the diversion effects---when I-710 is relieved, some cargo shifts to Oakland, potentially moving Oakland's connector from V/C 1.04 to V/C 1.25.

The port connector problem is solvable. The funding exists. The projects are well-understood. The binding constraint is classification, which determines funding eligibility, and coordination, which determines investment speed. Both are policy variables, not engineering ones.
